\section{Executive Summary}
\label{sec:executive-summary}

\subsection{Purpose and Audience}

This document sets out \textbf{governed pilot deployment requirements} for DUALEXIS---a
research-backed approach to \textbf{privacy-preserving situational awareness} and
\textbf{emergency coordination} in schools and similar confined educational spaces.

It is written for institutional decision-makers: university and school-group leadership,
safety coordinators, data protection officers, ethics committees, and Horizon Europe
consortium partners evaluating responsible innovation in educational safety.
The text supports structured planning and approval conversations; it is not legal advice,
regulatory certification, or evidence that a pilot will improve safety outcomes.

\subsection{What DUALEXIS Offers---and What It Does Not}

DUALEXIS is designed to help \textbf{licensed safety staff and administrators} build a
shared, zone-level picture of developing conditions---for example crowding near exits,
acoustic stress indicators, or conflicting signals in a hallway---and to receive
\textbf{advisory} recommendations that always pass through \textbf{human review} before
any institutional action.

Processing occurs \textbf{on school-controlled systems}.
Sensing is used only to produce short-lived, zone-scoped summaries; \textbf{video and
audio are not retained by default}, and the architecture excludes \textbf{facial
recognition, biometrics, and identity profiling}.

\begin{center}
\fbox{\parbox{0.92\textwidth}{\small
\textbf{Institutional boundaries (non-negotiable in pilot design).}
\begin{itemize}[noitemsep,leftmargin=*]
  \item \textbf{Advisory only} --- the system supports comprehension; it does not issue
        orders, trigger lockdowns, contact emergency services, or discipline students.
  \item \textbf{Human authority} --- protocol activation, public-address and radio use,
        parent communication, and escalation remain school responsibilities.
  \item \textbf{Privacy by design} --- no biometrics; no default media archive; purpose
        limited to safety coordination and accountable records.
  \item \textbf{Local control} --- institutions choose scope, retention, training, and
        whether any anonymized research export is enabled.
\end{itemize}
}}
\end{center}

DUALEXIS is \textbf{not} surveillance-as-usual, predictive policing, autonomous emergency
response, or a substitute for existing safety management systems (see Section~\ref{sec:what-dualxis-is-not}).
It does \textbf{not} imply GDPR compliance, EU AI Act conformity, or validated operational
effectiveness without site-specific legal review and ethics-approved evaluation.

\subsection{Why Institutions May Consider a Pilot}

Schools face a difficult balance: staff need timely awareness during stressful events,
yet communities rightly resist identity-centric monitoring and opaque automation.

A governed DUALEXIS pilot offers a \textbf{careful, deployable path} to explore whether
structured, privacy-preserving summaries can:

\begin{itemize}[noitemsep]
  \item strengthen \textbf{situational awareness} for safety teams without naming individuals;
  \item align with existing \textbf{lockdown, evacuation, and verification} procedures;
  \item document \textbf{accountability} through review queues and audit trails;
  \item generate reproducible evidence for research partnerships (including Horizon Europe)
        under explicit governance gates.
\end{itemize}

Success at this stage means \textbf{pilot readiness}: leadership authorization, DPIA
progress, protocol mapping, staff training, and ethics approval where required---not
proof that technology alone prevents harm.

\subsection{Recommended Pilot Pathway}

\begin{table}[htbp]
  \centering
  \caption{Indicative pilot pathway (planning guidance only).}
  \label{tab:pilot-tiers-summary}
  \small
  \begin{tabular}{@{}p{2.8cm}p{3.6cm}p{5.6cm}@{}}
    \toprule
    Phase & Scale & Institutional focus \\
    \midrule
    Foundation &
      1 school, limited zones &
      Governance charter, transparency plan, protocol mapping, tabletop exercises. \\
    Recommended pilot &
      1--3 schools &
      Staff training, bounded live sensing during staffed hours, workflow evaluation. \\
    Extended programme &
      5--10 schools &
      Multi-site learning, harmonized reporting, consortium replication (e.g.\ Horizon). \\
    \bottomrule
  \end{tabular}
\end{table}

Each phase should expand only after the prior governance gates are met.
Technical implementation detail for IT teams is consolidated in
Appendix~\ref{app:technical-architecture}; the main document emphasizes \textbf{roles,
accountability, and data-flow rules} institutions must approve.

\subsection{Governance Readiness Before Live Sensing}

No operational monitoring of occupants should begin until the steering committee confirms:

\begin{enumerate}[noitemsep]
  \item \textbf{Authorization} --- leadership approves scope, zones, and community communication.
  \item \textbf{Privacy and ethics} --- site-specific DPIA with qualified counsel; ethics
        committee sign-off where human-subjects evaluation is planned.
  \item \textbf{Protocol alignment} --- semantic event types mapped to the school's emergency
        plan (Section~\ref{sec:operational-protocols}).
  \item \textbf{Staff preparedness} --- training distinguishes \emph{advisory} system outputs
        from authoritative protocol steps; degraded-mode procedures documented.
  \item \textbf{Accountability} --- retention, access roles, false-alarm review, and pause/stop
        criteria agreed with safety and IT leads.
\end{enumerate}

\subsection{How to Use This Document}

The following sections translate these principles into objectives, school context, use cases,
privacy and DPIA checklists, operational workflows, evaluation planning, risks, timelines,
and indicative budgets.
Together they provide a \textbf{governance-first baseline} that research partners (UCJC/CRIA-BDHS)
and schools can adapt through controlled annexes---network plans, signed mappings, and
DPIA records---without weakening the safeguards above.

Institutions interested in pilot participation should begin with a steering-committee
workshop and DPO engagement.
Horizon partners should treat this document as a \textbf{responsible deployment template}:
privacy-preserving, human-supervised, locally controlled, and explicitly bounded in claims.
